\section{Anhang}

\subsection{Styleguide}

\begin{itemize}[itemsep=1pt, parsep=0pt]
    \item \textbf{Variablen, Konstanten}
    \begin{itemize}[itemsep=1pt, parsep=0pt]
        \item mit Kleinbuchstaben beginnen
        \item erster Buchstaben von zusammengesetzten Wörtern ist gross (mixed case)
        \item keine Underscores
    \end{itemize}
    Beispiele : counter, maxSpeed
    \item \textbf{Funktionen}
    \begin{itemize}[itemsep=1pt, parsep=0pt]
        \item mit Kleinbuchstaben beginnen
        \item erster Buchstaben von zusammengesetzten Wörtern ist gross (mixed case)
        \item Namen beschreiben Tätigkeiten
        \item keine Underscores
    \end{itemize}
    Beispiele : getCount(), init(), setMaxSpeed()
    \item \textbf{Klassen, Strukturen, Enums}
    \begin{itemize}[itemsep=1pt, parsep=0pt]
        \item mit Grossbuchstaben beginnen
        \item erster Buchstaben von zusammengesetzten Wörtern ist gross (mixed case)
        \item keine Underscores
        \item Namen beschreiben Dinge
    \end{itemize}
    Beispiele : MotorController, Queue, Color
\end{itemize}

\subsection{Code beispiele}
\subsubsection{String formatting}
Utilizza gli stream di output per costruire stringhe in memoria. È sicuro perché gestisce la memoria dinamicamente.

\begin{lstlisting}
#include <sstream>
#include <iomanip>

std::string formatTime(int h, int m) {
    std::ostringstream buf;
    buf << std::setfill('0') << std::setw(2) << h 
        << ":" << std::setw(2) << m;
    return buf.str();
}
\end{lstlisting}

Rappresenta lo standard moderno. Combina la sintassi concisa di \texttt{printf} con la sicurezza degli stream.

\begin{lstlisting}
#include <format>

std::string formatTime(int h, int m) {
    // :02 significa 2 cifre con riempimento zero
    return std::format("{:02}:{:02}", h, m);
}
\end{lstlisting}

\subsubsection{Ctor un dtor}
\lstinputlisting[language=c++]{code/ctor_dtor_anhang.cpp}

\subsection{Ritorno per Riferimento e Method Chaining}
\begin{itemize}
    \item \textbf{Meccanismo:} Una funzione restituisce un riferimento all'oggetto stesso
    utilizzando \texttt{return *this;}. Il tipo di ritorno deve 
    essere una referenza (es. \texttt{Complex\&}).
    \item \textbf{Kaskadierung (Cascading):} Permette di chiamare più metodi in 
    sequenza sulla stessa istanza in un'unica riga: \texttt{obj.funz1().funz2();}.
    
    \crd{Tramite ritorno \textit{by-value} funz2 verrebbe eseguita su un oggetto \textit{tmp}}
    \item \textbf{Efficienza:} A differenza del ritorno per valore, 
    non viene creata alcuna copia dei dati (nessun oggetto temporaneo \textit{tmp}), 
    risparmiando risorse 
    \item \textbf{Persistenza:} Le modifiche effettuate nella catena agiscono 
    direttamente sull'oggetto originale e non su una sua copia destinata alla distruzione.
\end{itemize}

\lstinputlisting[language=c++]{code/method_chaining.cpp}